\documentclass[11pt,letterpaper]{article}
\usepackage[margin=0.93in]{geometry}
\usepackage[T1]{fontenc}
\usepackage{lmodern,amsmath,amssymb,amsthm,microtype,xcolor}
\usepackage[colorlinks=true,allcolors=blue!48!black,pdfauthor={Blake Stone},pdftitle={Tilings by one bar of every length}]{hyperref}
\newtheorem{theorem}{Theorem}
\newtheorem{proposition}[theorem]{Proposition}
\raggedbottom
\begin{document}
\enlargethispage{2\baselineskip}
\begin{center}
{\LARGE\bfseries Tilings by one bar of every length}\par\vspace{7pt}
{\large A four-to-one rule for OEIS A394981}\par\vspace{9pt}
Blake Stone\par{\small Proof developed with OpenAI Codex}
\end{center}

In \href{https://oeis.org/A394981}{OEIS A394981}, Gary Yane asks about tilings by one $1\times j$ bar for each $1\le j\le n$. The bars tile a minimum-perimeter integer rectangle that accommodates them. Rotations of a bar are allowed, and tilings related by a symmetry of the board represent one term in the count $a(n)$. We work first with tilings on a fixed board, so placements at opposite edges remain distinguishable.

\begin{theorem} For every $h\ge2$, $a(2h)=4a(2h-1)$. The remaining odd starting index has $a(1)=a(2)=1$.\end{theorem}

\noindent\textbf{The rows.}
The boards for $n=2h-1$ and $n=2h$ have height $h$ and respective widths $2h-1$ and $2h+1$. Indeed, the longer side must be at least $n$, so the shorter side is at most $\lfloor(n+1)/2\rfloor$. This height minimizes perimeter at fixed area and is attainable. On the odd board use one full row and pair the remaining lengths to sum to $2h-1$. On the even board pair complementary lengths to sum to $2h+1$.

On the odd board the length-$2h-1$ bar fills a row. The length-$h$ bar therefore cannot run vertically, since it would have to cross that row. The $h$ bars of lengths $h,\ldots,2h-1$ are all horizontal. Each takes more than half the row, so there is exactly one in every row. On the even board the same holds for lengths $h+1,\ldots,2h$ by the height alone. Call these bars the long bars. Every long bar covers the middle column. Write the \emph{slack} of a row for the total number of cells to the left and right of its long bar. The row slacks, in some order, are
\[
    0,1,\ldots,h-1 \quad\text{on the odd board},\qquad
    1,2,\ldots,h   \quad\text{on the even board}.
\]
Every short bar lies entirely to one side of the middle column. Consequently we can lengthen all long bars by a common amount while fixing the short bars' distances to their respective outer edges. This amounts to inserting columns within the band occupied by the long bars. It preserves all incidences and vertical alignments on each side.

\begin{proof}[Proof of the theorem]
Fix an even-board tiling and inspect the short bar of length $h$.

\emph{Vertical.} This bar spans every row. In the row of slack one it occupies the unique cell left free by the length-$2h$ bar. That cell is on the left or right outer edge, so the vertical bar forms a complete outer column. Remove it. Now shorten each long bar by one, keeping the two edge patterns fixed. This gives the odd board, with one long bar of each length $h,\ldots,2h-1$ and all the short bars $1,\ldots,h-1$. Conversely, lengthen each odd-board long bar by one and append a full column on the chosen side. Each odd-board tiling has exactly two even-board preimages of this kind.

\emph{Horizontal.} The bar needs $h$ consecutive slack cells, so it lies in the unique row of slack $h$, entirely on one side of the length-$(h+1)$ bar. These two bars fill that row. Replace the row by a single full bar of length $2h-1$. In all other rows shorten each long bar by two, retaining the two edge patterns. Their lengths become $h,\ldots,2h-2$. The replaced row contained no vertical tiles. The inverse splits the full row of the odd board into lengths $h$ and $h+1$ in either order and lengthens the other long bars by two. This gives precisely two more preimages.

The two cases are disjoint for $h\ge2$. They exhaust all possibilities and give a factor of four for fixed-board counts. Both boards are nonsquare and their four symmetries act freely: a symmetry reversing row order cannot fix every row's unique long bar, and a left-right reflection cannot center the long bar in the row of slack one. Dividing both counts by four proves the assertion. For $h=1$ the length-one tile has no distinct orientations, and the $1\times1$ and $1\times3$ boards each give one symmetry class.
\end{proof}


\clearpage
\begin{center}
\begingroup\setlength{\unitlength}{0.135in}
\newcommand{\barbox}[5]{\put(#1,#2){\framebox(#3,#4){\scriptsize #5}}}
\newcommand{\highlightbar}[1]{\textcolor{blue!55!black}{\bfseries #1}}
\begin{tabular}{c@{\quad}c@{\quad}c@{\quad}c@{\quad}c}
\begin{picture}(3,2)\barbox{0}{1}{3}{1}{3}\barbox{0}{0}{2}{1}{2}\barbox{2}{0}{1}{1}{1}\end{picture}&
\begin{picture}(5,2)\barbox{0}{0}{1}{2}{\highlightbar{2}}\barbox{1}{1}{4}{1}{4}\barbox{1}{0}{3}{1}{3}\barbox{4}{0}{1}{1}{1}\end{picture}&
\begin{picture}(5,2)\barbox{4}{0}{1}{2}{\highlightbar{2}}\barbox{0}{1}{4}{1}{4}\barbox{0}{0}{3}{1}{3}\barbox{3}{0}{1}{1}{1}\end{picture}&
\begin{picture}(5,2)\barbox{0}{1}{2}{1}{\highlightbar{2}}\barbox{2}{1}{3}{1}{3}\barbox{0}{0}{4}{1}{4}\barbox{4}{0}{1}{1}{1}\end{picture}&
\begin{picture}(5,2)\barbox{0}{1}{3}{1}{3}\barbox{3}{1}{2}{1}{\highlightbar{2}}\barbox{0}{0}{4}{1}{4}\barbox{4}{0}{1}{1}{1}\end{picture}\\
\scriptsize Odd board&\scriptsize Vertical, left&\scriptsize Vertical, right&\scriptsize Horizontal, left&\scriptsize Horizontal, right
\end{tabular}
\endgroup

\small One $2\times3$ tiling and its four $2\times5$ extensions. Labels give bar lengths.
\end{center}
\section*{Counting every term}
A horizontal grid line is a \emph{cut} if it crosses no tile interior. Count both horizontal boundary lines. Let $E_h(u)$ and $O_h(u)$ be the polynomials that count fixed-board tilings for $2h$ and $2h-1$ respectively, assigning weight $u^c$ to a tiling with $c$ cuts.

\begin{proposition} We have $E_1(u)=2u^2$, $O_1(u)=u^2$ and, for $h\ge2$,
\[
 O_h(u)=u^2 E'_{h-1}(u),\qquad
 E_h(u)=2O_h(u)+2O_h(1)u^2 .
\]
In particular $a(2h-1)=O_h(1)/4$ and $a(2h)=E_h(1)/4$ for $h\ge2$.
\end{proposition}
\begin{proof}
The even board of height $h-1$ and the odd board of height $h$ have the same width $2h-1$. Insert the full row at any cut of the smaller board. Removing the unique full row is the inverse. An object with $c$ cuts gives $c$ larger objects, each with $c+1$ cuts. In the four constructions above, the two horizontal choices preserve every cut. The two vertical choices leave only the two boundary cuts.
\end{proof}

For example, $O_3(u)=16u^3+24u^4$, giving $a(5)=10$. The next terms after the twenty listed in OEIS revision 42 are
\[
\begin{array}{c|rrrr}
n&21&22&23&24\\\hline
a(n)&21{,}459{,}460{,}096&85{,}837{,}840{,}384&464{,}513{,}597{,}440&1{,}858{,}054{,}389{,}760 .
\end{array}
\]

\section*{Growth}
Let $e_h=E_h(1)$ and $t_h=e_h/(2^h h!)$. For $h\ge2$ the recurrence reads
$E_h=2u^2 E'_{h-1}+(e_h/2)u^2$. The operator $D=2u^2(d/du)$ sends $u^2$ after $k$ iterations to $2^k(k+1)!u^{k+2}$. Iterating from $E_1=2u^2$, evaluating at one and moving the $e_h/2$ term to the left gives
\[
 t_h=2+\sum_{j=2}^{h-1}w_{h,j}t_j,\qquad
 w_{h,j}=\frac{j!(h-j+1)!}{h!}=\frac{h+1}{\binom{h+1}{j}} .
\]
The end weights $w_{h,2}$ and $w_{h,h-1}$ both equal $2/h$ for $h\ge4$. The remaining weights sum to $O(h^{-2})$: the terms at $j=3,h-2$ are each $O(h^{-2})$, and each of the $O(h)$ terms at $4\le j\le h-3$ is at most $(h+1)/\binom{h+1}{4}$. Hence the full weight sum is $O(1/h)$. Once this sum is at most $1/2$, induction bounds every $t_h$ by the larger of four and the maximum of the finitely many preceding values. Substitution gives $t_h=2+O(1/h)$. Finally $t_2=2$; substituting the two end terms gives $t_h=2+8/h+O(h^{-2})$. Thus
\[
 \boxed{\displaystyle a(2h)=2^{h-1}h!\left(1+\frac4h+O(h^{-2})\right).}
\]
\end{document}
